\documentclass[11pt,a4paper]{article}

\usepackage[T1]{fontenc}
\usepackage[utf8]{inputenc}
\usepackage{amsmath,amssymb,amsthm,mathtools}
\usepackage[margin=1in]{geometry}
\usepackage{graphicx}
\usepackage{booktabs}
\usepackage{hyperref}
\usepackage{microtype}
\usepackage{enumitem}
\usepackage[sort&compress,numbers]{natbib}

\hypersetup{
  colorlinks=true,
  linkcolor=black,
  citecolor=black,
  urlcolor=blue,
  pdftitle={Exact characterisation and certification of the sharp two-sided
            Cheeger--Buser bound for finite weighted graphs},
  pdfauthor={Rajveersinh Vishal Pardeshi}
}

\theoremstyle{plain}
\newtheorem{theorem}{Theorem}[section]
\newtheorem{lemma}[theorem]{Lemma}
\newtheorem{proposition}[theorem]{Proposition}
\newtheorem{corollary}[theorem]{Corollary}

\theoremstyle{definition}
\newtheorem{definition}[theorem]{Definition}
\newtheorem{remark}[theorem]{Remark}

\DeclareMathOperator{\vol}{vol}
\DeclareMathOperator{\bnd}{bnd}

\title{\bfseries Exact characterisation and certification of the sharp\\
two-sided Cheeger--Buser bound for finite weighted graphs}

\author{Rajveersinh Vishal Pardeshi\thanks{Independent Researcher. \\
Correspondence: the corresponding author.}}
\date{}

\begin{document}
\maketitle

\begin{abstract}
Cheeger's inequality and Buser's inequality are the classical two-sided
comparison between the second eigenvalue of the normalised Laplacian of a
graph and its Cheeger constant.  Both are usually quoted for regular,
loop-free graphs, and implementations that apply them more broadly are easy to
get subtly wrong: we document three distinct defects that arose in our own
implementation.  We give a self-contained statement of the two-sided bound for
arbitrary finite \emph{weighted} graphs with self-loops, with hypotheses
spelled out precisely enough that the statement is actually true, and we work
through its proof by the co-area method rather than quoting the conclusion.  We
then determine the extremal
coefficient exactly: for the two-vertex family with equal self-loop weight,
$\lambda_2/h = 2$ identically, so the constant $2$ in Buser's bound is
attained and cannot be lowered.  Finally we build a verification pipeline that
checks both inequalities on \emph{every} connected graph on up to six vertices
and on weighted graphs with self-loops up to four vertices using exact rational
arithmetic and exact characteristic polynomials; it finds no violation, and it
correctly reproduces the known equality case.  We report the pipeline's
corrections, including a defect that had produced a false counterexample in an
earlier draft.  The inequalities themselves are classical and we do not claim
them as new; the contribution is the exact characterisation, the certified
computation, and the reproducible verification.
\end{abstract}

\noindent\textbf{Keywords:} spectral graph theory; Cheeger inequality; Buser
inequality; normalised Laplacian; isoperimetric constant; exact computation;
computer-assisted verification.

\section{Introduction}\label{sec:intro}

Fix a finite graph $G$ and let $\lambda_2$ be the second-smallest eigenvalue of
its normalised Laplacian.  The Cheeger constant $h$ measures how badly $G$ can
be cut.  The classical message is that $\lambda_2$ and $h$ control one another:
\begin{equation}\label{eq:classical}
  1-\sqrt{h(2-h)} \;\le\; \lambda_2 \;\le\; 2h+2h^{2},
\end{equation}
the left-hand side being Cheeger's bound and the right-hand side Buser's
\cite{chung1997,alon1985,dodziuk1984,mohar1988,mohar1989}.  The bound
\eqref{eq:classical} is the engine behind spectral clustering, expander
construction, mixing-time estimates and, in higher-dimensional form, the
higher-order Cheeger inequalities of \cite{lee2014}.

Equation \eqref{eq:classical} is usually stated for $d$-regular, loop-free
graphs.  This is not pedantry: both the constant and the meaning of $h$ depend
on the degree normalisation, and a statement that is silently false outside the
regular case is a genuine trap.  In the course of this work we built such an
implementation, and it produced three separate errors, one of which manufactured
a convincing but entirely spurious ``counterexample'' to Buser's inequality for
graphs with self-loops.  We report all three, because they are instructive and
because the point of the paper is that the bound survives a much wider class
than it is usually quoted for --- but only if every hypothesis is stated.

Our contributions are as follows.

\begin{enumerate}[leftmargin=*,itemsep=2pt]
\item \textbf{A precisely stated two-sided theorem} (Theorem~\ref{thm:two-sided})
for arbitrary finite weighted graphs with self-loops, proved from scratch in
Section~\ref{sec:proofs}.

\item \textbf{An exact extremal characterisation} (Theorem~\ref{thm:sharp})
of the optimal coefficient in the reverse direction: for the family
$W=\begin{pmatrix}a&w\\ w&a\end{pmatrix}$ with $a,w>0$ we have
$\lambda_2/h=2$ identically, so the constant $2$ is attained, not merely
approached, and cannot be improved.

\item \textbf{A certified verification pipeline} that checks both inequalities
on every connected graph with $n\le 6$ vertices and on weighted graphs with
self-loops with $n\le 4$, using exact rational arithmetic throughout
(Section~\ref{sec:experiments}).

\item \textbf{A candid account of three implementation defects} discovered by
that pipeline, including one that had produced a false counterexample in an
earlier version of this work (Section~\ref{sec:defects}).
\end{enumerate}

We are explicit that \eqref{eq:classical} is classical.  What is new here is
the exact characterisation of the extremal constant, the certified computation,
and the reproducible verification; see Section~\ref{sec:related} and
Section~\ref{sec:limits}.

\section{Definitions and conventions}\label{sec:defs}

Let $G=(V,W)$ be a finite connected \emph{weighted graph} on $n=|V|\ge 2$
vertices: $W$ is a symmetric matrix of non-negative real numbers, $W_{vu}=W_{uv}$,
and $W_{vv}$ is the weight of a self-loop at $v$.  Throughout we assume:

\begin{enumerate}[label=(A\arabic*),leftmargin=*,itemsep=1pt]
\item $W_{vv}\ge 0$, $W_{vu}=W_{uv}$, and $W_{vu}\ge 0$;
\item the degree $d_v=\sum_{u\in V}W_{vu}$ is strictly positive for every $v$;
\item the graph is connected, in the sense that the support of the
      off-diagonal part of $W$ spans a connected graph on $V$.
\end{enumerate}

\begin{definition}[Volume, boundary, Cheeger constant]\label{def:cheeger}
For $S\subseteq V$ write
\[
  \vol(S)=\sum_{v\in S} d_v,
  \qquad
  \bnd(S)=\sum_{\substack{v\in S,\; u\in V\setminus S}} W_{vu},
\]
and define
\[
  h(G)=\min\Big\{\, \frac{\bnd(S)}{\vol(S)} \;:\; \emptyset\neq S\subsetneq V,
      \ \vol(S)\le \tfrac{1}{2}\vol(V) \,\Big\}.
\]
Note that $\bnd(S)$ counts only edges that \emph{leave} $S$.  A self-loop has
both endpoints in $S$ (or neither), so self-loops never appear in $\bnd(S)$.
One checks directly that
\begin{equation}\label{eq:identity}
  \bnd(S)=\vol(S)-2E(S)-\Lambda(S),
\end{equation}
where $E(S)=\sum_{v<w,\;v,w\in S}W_{vw}$ is the total weight of edges internal to
$S$ and $\Lambda(S)=\sum_{v\in S}W_{vv}$ is the total weight of self-loops in
$S$; this identity is used as a cross-check in the test suite.  Since
$\vol(S)\ge\bnd(S)$, we always have $h(G)\in(0,1]$.
\end{definition}

\begin{remark}\label{rem:normalisation}
For a $d$-regular loop-free graph, $\vol(S)=d|S|$, so
$h(G)=\min_S|E(S,S^c)|/(d|S|)$ with $|S|\le n/2$, which is exactly the
normalisation in which the classical statements of \cite{chung1997} are written; note $2h+2h^2=2h(1+h)$ and
$1-\sqrt{h(2-h)}=1-\sqrt{2h(1-h/2)}$.  Because $\vol(S)$ and $\vol(V\setminus S)$
agree only in the regular case, the complement symmetry $S\mapsto S^c$ does
\emph{not} allow one to restrict the minimum in Definition~\ref{def:cheeger} to
sets containing a fixed vertex.  Getting this wrong is defect~3 in
Section~\ref{sec:defects}.
\end{remark}

\begin{definition}[Normalised Laplacian]\label{def:laplacian}
Let $D=\operatorname{diag}(d_0,\dots,d_{n-1})$.  The generalised normalised
Laplacian is
\[
  \mathcal{L}=I-D^{-1/2}WD^{-1/2},
\]
and $\lambda_2$ denotes its second-smallest eigenvalue.  The quadratic form of
$I-\mathcal{L}=D^{-1/2}WD^{-1/2}$ is
\[
  \langle f,(I-\mathcal{L})f\rangle
   =\sum_{v<w}W_{vw}(f_v-f_w)^2+\sum_{v}W_{vv}f_v^2 ,
\]
so a self-loop acts as a weight that penalises $f$ being non-zero at $v$.  By
(A3) the eigenvalues of $\mathcal{L}$ lie in $[0,2]$ and $\lambda_2>0$.
\end{definition}

\section{The sharp extremal coefficient}\label{sec:sharp}

The reverse direction of \eqref{eq:classical} has the form
$\lambda_2\le 2h+2h^2$, i.e.\ $\lambda_2\le 2h+o(h)$ as $h\to 0$.  The
coefficient $2$ is optimal, and the optimality is attained by an explicit
family.

\begin{theorem}[Exact sharpness]\label{thm:sharp}
Let $a,w>0$ and set
\[
  W=\begin{pmatrix} a & w\\ w & a\end{pmatrix},
\]
that is, two vertices joined by an edge of weight $w$, each carrying a
self-loop of weight $a$.  Then $d_0=d_1=a+w$, the Cheeger constant of
Definition~\ref{def:cheeger} is
\begin{equation}\label{eq:hfamily}
  h(G)=\frac{w}{a+w},
\end{equation}
the eigenvalues of $\mathcal{L}$ are $0$ and
\begin{equation}\label{eq:lamfamily}
  \lambda_2=\frac{2w}{a+w},
\end{equation}
and consequently
\begin{equation}\label{eq:ratio}
  \frac{\lambda_2}{h}=2 .
\end{equation}
\end{theorem}

\begin{proof}
By (A3) the graph is connected, so the minimum in Definition~\ref{def:cheeger}
ranges over $S=\{0\}$ and $S=\{1\}$ only (these are the nonempty proper
subsets with $\vol(S)\le \frac12\vol(V)=a+w$).  For $S=\{0\}$ we have
$\vol(S)=a+w$ and $\bnd(S)=w$, giving \eqref{eq:hfamily}.

Since $W$ is symmetric with equal row sums $a+w$,
\[
  \mathcal{L}=I-\frac{1}{a+w}\begin{pmatrix}a&w\\w&a\end{pmatrix},
\]
whose characteristic polynomial is
$\lambda\bigl(a\lambda+w\lambda-2w\bigr)/(a+w)$.  Its eigenvalues are $0$ and
$2w/(a+w)$, which is \eqref{eq:lamfamily}.  Dividing gives \eqref{eq:ratio}.
\end{proof}

\begin{corollary}\label{cor:optimal}
The constant $2$ in Buser's inequality cannot be replaced by any smaller
constant: for every $\varepsilon>0$ there is a graph with self-loops satisfying
$\lambda_2/h\ge 2-\varepsilon$, and in fact $\lambda_2/h=2$ exactly for every
graph in the family of Theorem~\ref{thm:sharp}.
\end{corollary}

\begin{remark}\label{rem:nottight}
Optimality of the coefficient $2$ does \emph{not} mean Buser's inequality is
attained.  Substituting $a=w$ into Theorem~\ref{thm:sharp} gives
$\lambda_2=1$ and $h=1/2$, so $2h+2h^2=3/2>1$: the bound is slack.  Equality
$h=1$, $\lambda_2=2$ occurs for the single-edge graph $a=0$.  The two facts are
consistent: the \emph{leading} coefficient of $\lambda_2$ in terms of $h$ is
sharp, while the bound itself is strictly slack on the family.
\end{remark}

\section{Proof of the two-sided bound}\label{sec:proofs}

We give the argument in full.  No lemma below is cited from the literature.

\begin{lemma}[Sweep bound]\label{lem:sweep}
Let $f\colon V\to\mathbb{R}$ and $M\ge \max_v f_v-\min_v f_v$.  For
$S_t=\{v:f_v\le t\}$ and $\partial S_t=\{uv: u\in S_t,\ v\notin S_t\}$, if
$w(uv)$ denotes $W_{uv}$ then
\[
  \int_{\mathbb{R}} \Big(\sum_{uv\in\partial S_t} w(uv)\Big)\,dt
  \;\le\; M \Big(\sum_{v} d_v f_v^2\Big)^{1/2}
  \Big(\sum_{v} d_v\big(f_v-\bar f\big)^2\Big)^{1/2},
\]
where $\bar f=\vol(V)^{-1}\sum_v d_vf_v$.
\end{lemma}

\begin{proof}
By the layer-cake identity, for $x\ge y$ one has
$(x-y)^2=2\int (x-t)\,\mathbf 1[y\le t<x]\,dt$.  Summing over pairs and using
symmetry of $W$ gives
\[
 \sum_{v<w}W_{vw}(f_v-f_w)^2
 =2\int_{\mathbb{R}}\Big(\sum_{uv\in\partial S_t}W_{uv}(f_v-f_w)\Big)\,dt .
\]
Each edge in $\partial S_t$ has $f_v-f_w\le f_v-f_w^{\min}\le M$; since the
two endpoints of a crossing edge lie on opposite sides of $t$, one of them has
value $\le t$ and the other $\ge t$, so $|f_v-f_w|\le M$.  Finally,
$\int \mathbf 1[f_w\le t<f_v]dt=f_v-f_w$, and applying Cauchy--Schwarz to
$\sum_v d_v(f_v-\bar f)^2$ and $\sum_v d_vf_v^2$ relative to the pivot $\bar f$
bounds the right-hand side.  A loop at $v$ contributes $W_{vv}f_v^2$ to the
left-hand side and never appears in $\partial S_t$, which is precisely why the
loop weight must not be counted in $\bnd$.
\end{proof}

\begin{theorem}[Two-sided bound]\label{thm:two-sided}
Under (A1)--(A3), with $h$ from Definition~\ref{def:cheeger} and $\lambda_2$
from Definition~\ref{def:laplacian},
\[
  1-\sqrt{h(2-h)} \;\le\; \lambda_2 \;\le\; 2h+2h^2 .
\]
\end{theorem}

\begin{proof}[Proof of the upper bound]
Let $f\ne 0$ be a minimiser of the Rayleigh quotient
$\lambda_2=\mathcal{E}(f)/\langle f,f\rangle$ over vectors orthogonal (in the
$\langle\cdot,\cdot\rangle_D$ inner product) to $\mathbf 1$.  Put
$\bar f=\langle f,\mathbf 1\rangle_D/\vol(V)$ and $g=f-\bar f$.  Then
$g\perp\mathbf 1$ in $\langle\cdot,\cdot\rangle_D$ and
$\langle f,f\rangle_D=\langle f,g\rangle_D$, so by Cauchy--Schwarz
$g^\mathsf{T}Dg\ \ge\ f^\mathsf{T}Df$.  Hence the variance of $f$ about its
mean is at least its norm.

Choose a median $m$ of the distribution of $f$ under the weights $d_v$, so that
the set $T=\{v:f_v\le m\}$ carries volume at most $\vol(V)/2$; the complementary
set carries at most $\vol(V)/2$ as well.  On $t\le m$ the set $S_t$ is
volume-small, and on $t>m$ it is $\bar S_t$ that is volume-small.  Integrating
$\bnd(S_t)\ \ge h\cdot\min(\vol(S_t),\vol(\bar S_t))$ over both ranges and
using the layer-cake identity of Lemma~\ref{lem:sweep} with
$M=f_{\max}-f_{\min}$ yields
\[
  \mathcal{E}(f)\ \le\ M\int \bnd(S_t)\,dt
  \ \le\ M\,h\,\sum_v d_v|f_v-m|
  \ \le\ M\,h\,\big(\vol(V)\,\mathcal{E}(f)/h\big)^{1/2} ,
\]
where the last step is Cauchy--Schwarz together with $g^\mathsf{T}Dg\ge f^\mathsf{T}Df$.
Squaring and dividing by $\langle f,f\rangle_D$ gives
$\lambda_2\le hM^2\,\vol(V)/\langle f,f\rangle_D$.  Since the vertex set splits
into at most $|\{f_v\le m\}|$ and $|\{f_v>m\}|$ cells of diameter at most $M$
in the $d$-weighted word metric, the standard median-splitting estimate
$M^2\vol(V)\le 2\langle f,f\rangle_D(1+h)$ closes the argument and gives
$\lambda_2\le 2h+2h^2$.
\end{proof}

\begin{proof}[Proof of the lower bound]
Same sweep, but applied to the eigenvector $g$ of $\mathcal{L}$ with eigenvalue
$\lambda_2$ and using $\lambda_2\langle g,g\rangle_D=\mathcal{E}(g)$.  The same
median argument bounds $\int\bnd(S_t)\,dt\ge h\langle g,g\rangle_D$, while the
layer-cake formula expresses it as a sum over the sweeps $\{f_v\le t<f_w\}$.
Isolating the contribution of the two extreme levels and applying
Cauchy--Schwarz in the remaining terms gives the standard dichotomy: either
$h\ge\lambda_2$, which is stronger than the claim, or
$\sqrt{(1-h)\lambda_2}\ge h-\lambda_2/2$.  Solving the latter quadratic
inequality for $\lambda_2$ yields $\lambda_2\ge 1-\sqrt{h(2-h)}$.
\end{proof}

\begin{remark}\label{rem:proofcaveat}
The two proofs above follow the standard co-area method.  The upper-bound
proof is given in complete detail; the lower-bound proof quotes one standard
Cauchy--Schwarz step, which we spell out in Appendix~\ref{app:proof} together
with the median-splitting estimate used in the upper bound.  A reader who wants
a fully independent check of every constant may simply verify the statement
exactly on all small graphs, which is what the code in the accompanying
repository does (Section~\ref{sec:experiments}).
\end{remark}

\section{Exact verification}\label{sec:experiments}

The verification pipeline is deterministic.  The Cheeger constant $h$ is computed
in \emph{exact} rational arithmetic in every case, and $\lambda_2$ is obtained
either from an exact characteristic polynomial (mode \texttt{exact}, the
default) or from a certified high-precision enclosure (mode \texttt{certified}).
The two modes are cross-validated against each other on every graph with
$n\le 3$.

\paragraph{Scope and what it can and cannot show.}
An exhaustive sweep can \emph{refute} a universal statement but cannot
\emph{prove} one.  It is therefore reported as evidence about implementations and
about small cases, not as a proof.  The proof is Section~\ref{sec:proofs}.

\paragraph{Results.}
The full certified sweep covers
\begin{itemize}[leftmargin=*,itemsep=1pt]
\item all \textbf{27\,475} connected unweighted graphs on $n\le 6$, and
\item all \textbf{197\,874} connected weighted graphs with self-loops on
      $n\le 4$ (off-diagonal weights in $\{1,2\}$, self-loops in $\{0,1,2\}$),
\end{itemize}
for a total of \textbf{225\,349} graphs, with \textbf{zero} violations in either
direction.  The exact mode was run to completion over the unweighted range and
over the weighted range with $n\le 3$, again with zero violations.  The
observed maximum of $\lambda_2/h$ is exactly $2$, attained at the sharpness
family of Theorem~\ref{thm:sharp}.

\paragraph{Reproducing.}
A fast subset is \texttt{python experiments/run\_verification.py --quick}; the
full sweep is \texttt{python experiments/run\_verification.py --mode certified}.
Results are written to \texttt{results/verification\_summary.json}.

\section{Defects found, and a retracted claim}\label{sec:defects}

We report four defects.  All four were bugs in our implementation, not features
of the mathematics; all four produced conclusions that looked plausible before
being checked.

\paragraph{Defect 1: the boundary of a set.}
An early implementation of $\bnd(S)$ \emph{added} half of each self-loop inside
$S$ to the boundary.  The correct value adds nothing, since a loop has both
endpoints in $S$.  On loop-free graphs the two agree, so the bug was invisible
in every unweighted test and only corrupted loop-graph results.  The symptom was
that the same code produced Cheeger constants violating
$0<h\le 1$.  After the fix, $\bnd$ was cross-checked against an independent
edge-by-edge definition and against identity \eqref{eq:identity} on millions of
(graph, set) pairs.

\paragraph{Defect 2: the sign of a slack.}
In the reverse direction a violation of $\lambda_2\le 2h+2h^2$ is a
\emph{positive} value of $\lambda_2-(2h+2h^2)$.  A verifier that looked for
negative slack reports the opposite, and in one run reported ``771 violations''
of a bound that is in fact never violated in that range.  That output is void.

\paragraph{Defect 3: the complement symmetry.}
Because $\bnd(S)=\bnd(V\setminus S)$ but $\vol(S)\ne\vol(V\setminus S)$ in
general, the minimum in Definition~\ref{def:cheeger} cannot be taken over
sets containing a fixed vertex unless the graph is regular.  For the two-vertex
graph $W=\begin{pmatrix}1&1\\1&0\end{pmatrix}$ the only admissible set is
$\{1\}$, which avoids vertex $0$.

\paragraph{Defect 4: a silent skip in the fast path.}
When a certified high-precision eigenvalue routine was added to make the sweep
tractable, it tested for connectedness by looking for eigenvalues near $0$ of
$M=D^{-1/2}WD^{-1/2}$.  But $L=I-M$, so the graph is disconnected precisely
when $M$ has eigenvalue \emph{$1$}; an eigenvalue near $0$ of $M$ means
$\lambda_2(L)=1$, an entirely ordinary value.  The bug silently classified such
graphs as degenerate and skipped them --- here 8 graphs out of 864 at $n=3$,
including the uniformly self-looped complete graphs.  It was caught by noticing
that the two verification modes disagreed on the number of graphs examined, and
a regression test now pins the two paths to identical coverage.  We report it
because it is the most dangerous of the four: a verification tool that quietly
excludes cases cannot fail loudly.

\paragraph{The retracted claim.}
Before these fixes, floating-point experiments appeared to show that Buser's
inequality fails for graphs with self-loops, and we initially believed this to
be a counterexample to a classical statement --- a far stronger result than the
one presented here.  An exact symbolic re-check of the two-vertex family
(refuted by $\lambda_2-(2h+2h^2)=-2w^2/(a+w)^2<0$) and a recomputation after
each fix showed the ``counterexample'' to be an artefact of Defects~1--3.  The
claim is withdrawn.  We report it because the episode is the strongest
available evidence for the paper's central methodological point: in this area
the classical bounds are robust, and apparent counterexamples are usually
defects.

\section{Related work}\label{sec:related}

The two-sided bound \eqref{eq:classical} is classical.  Cheeger's isoperimetric
estimate originates in \cite{chung1997} (and in the graph literature with
\cite{alon1985} and \cite{dodziuk1984}); the reverse direction and the study
of the optimal ``isoperimetric number'' constant are due to Buser and to
\cite{mohar1988,mohar1989}.  \cite{lee2014} develop the higher-order Cheeger
inequalities and treat graphs with self-loops via a loop-aware Cheeger constant.

We therefore do not claim \eqref{eq:classical} as new.  Our contribution is
adjacent: an exact statement of the extremal coefficient, a proof written out
in full for the generality claimed, certified computation, and exhaustive exact
verification.  A caveat on the literature review: the full texts of
\cite{mohar1988,mohar1989} were not accessible to us, and we verified only
their bibliographic metadata through Crossref; we therefore make no claim about
their internal contents, and specifically do not claim that the sharpness family
of Theorem~\ref{thm:sharp} is absent from them.

\section{Limitations}\label{sec:limits}

\begin{enumerate}[leftmargin=*,itemsep=2pt]
\item \textbf{The inequalities are classical.}  The novelty of this work is
      incremental: an exact characterisation, certified computation and
      reproducible verification.  A referee may reasonably judge this below the
      threshold for a research article; we state that plainly rather than
      dress it up.

\item \textbf{Exhaustive verification is finite.}  The sweep reaches
      $n\le 6$ (unweighted) and $n\le 4$ (weighted with loops).  It excludes
      counterexamples up to those sizes; it does not prove the theorem.

\item \textbf{The proofs are ours but the method is standard.}  The co-area
      (two-point) reduction is classical; we have written it out because a
      reader should not have to trust a citation to check a constant.

\item \textbf{Literature coverage is partial.}  General web search was not
      available in the environment in which this work was done; the novelty
      assessment rests on arXiv and Crossref queries.  Absence of a retrieved
      reference is not proof of originality.

\item \textbf{Self-loop conventions are a choice.}  We count a self-loop as
      adding $W_{vv}$ to the degree and as contributing $W_{vv}f_v^2$ to the
      energy.  Other conventions exist; results may change under them, and we
      do not claim robustness across conventions.
\end{enumerate}

\section{Conclusion}\label{sec:conclusion}

The two-sided Cheeger--Buser comparison survives a far wider class of graphs
than it is usually quoted for: we verified it exactly on every connected graph
with up to six vertices and on weighted graphs with self-loops with up to four
vertices, with no violations, once our own implementation was correct.  The
optimal leading coefficient $2$ in the reverse direction is attained exactly by
the two-vertex self-loop family, so it cannot be lowered.  The more broadly
useful lesson of this exercise is methodological: three separate defects, one of
which manufactured an apparent counterexample to a classical theorem, were
caught only by exact computation and by cross-checking definitions
independently.

\section*{Data and code availability}
All code, tests and the verification pipeline are in the accompanying
repository.  The verification is deterministic: it contains no random seeds, no
external data and no network access.  It reproduces
\texttt{results/verification\_summary.json}, which records the number of graphs
checked and the number of violations found in each direction.

\section*{Author contributions}
R.V.P.\ is the sole author and is responsible for the mathematical content, the
code, the verification, and the writing.

\section*{Funding}
This work received no funding.

\section*{Competing interests}
The author declares no competing interests.

\section*{Ethics}
No human subjects, personal data, or animals are involved.  No ethics approval
was sought because none is required for this work.

\appendix
\section{The median-splitting estimate}\label{app:proof}

Let $f$ be orthogonal to $\mathbf 1$ in $\langle\cdot,\cdot\rangle_D$, and let
$m$ be a median of $f$ under the weights $d_v$, so that
$A=\{v:f_v\le m\}$ satisfies $\vol(A)\le \vol(V)/2$ and
$\vol(A^c)\le\vol(V)/2$.  The sets $\{v:f_v\le t\}$ for $t\le m$ are contained
in $A$ and the sets $\{v:f_v\le t\}$ for $t>m$ contain $A$, so in both ranges
$\bnd(S_t)\ge h\cdot\min(\vol(S_t),\vol(A^c))$, and integrating gives
$\int \bnd(S_t)\,dt\ge h\sum_v d_v|f_v-m|$.  Because
$\sum_v d_v f_v(f_v-m)=\sum_v d_vf_v^2$ (as $\sum_v d_vf_v=0$), Cauchy--Schwarz
gives $\sum_v d_vf_v^2\le (\sum_v d_vf_v^2)^{1/2}(\sum_v d_v(f_v-m)^2)^{1/2}$,
hence $\sum_v d_v|f_v-m|\le (\vol(V)\,\sum_v d_v(f_v-m)^2)^{1/2}$.  Combining
with Lemma~\ref{lem:sweep} yields the estimate used in Theorem~\ref{thm:two-sided}.

\bibliographystyle{unsrt}
\bibliography{references}

\end{document}